\section{Data and Experimental Protocol}\label{sec:protocol}
\subsection{Collections, splits, and model identities}
Table~\ref{tab:datasets} records the collections, complete judged test-query counts, and frozen base settings. Every positive-query ID in each official test split is evaluated: 300 SciFact, 323 NFCorpus, and 648 FiQA queries. We exclude the partially collected fourth collection. The separate resource study uses saved features from the same domains, not an additional retrieval-effectiveness test.

\begin{table*}[t]
\centering\small
\caption{Recorded datasets and configurations. Documents are parent documents. SciFact development is derived from official training; the other two use official development splits. The lexical weight is $\beta$ and the base redundancy weight is $\lambda$.}
\label{tab:datasets}\begin{tabular}{lrrrrrl}
\toprule
Dataset & Documents & Dev. & Test & $\beta$ & $\lambda$ & Reported policy\\
\midrule
SciFact & 5,183 & 161 & 300 & 0.25 & 0.000 & Base + refinement\\
NFCorpus & 3,633 & 324 & 323 & 0.50 & 0.050 & Base + refinement\\
FiQA & 57,638 & 500 & 648 & 0.00 & 0.025 & Base\\
\bottomrule
\end{tabular}

\end{table*}

The dense encoder, \path{sentence-transformers/all-MiniLM-L6-v2}, produces normalized 384-dimensional vectors. For relevance we use \path{cross-encoder/ms-marco-MiniLM-L6-v2}, with an L2 cross-encoder as the smaller neural control. Repository revisions and model-file hashes in the evidence manifest identify the artifacts used. The journal analysis downloads no new weights and reruns no inference.

We normalize document text with Unicode NFKC and whitespace collapse, using deterministic text handling. Tokenizer offsets are computed without padding or truncation. Chunks contain at most 190 WordPiece content tokens and advance by 158 content tokens; adding two special tokens gives a model input of at most 192 tokens. The maximum query cosine over a parent's chunks becomes its score, with the lowest chunk index breaking an exact tie. We use that original chunk as the representative passage, so methods sharing a pool also share parent IDs and passage representations.

Query lengths must satisfy the content-token limit, and query--passage pairs must satisfy the cross-encoder's 512-token limit; neither check permits silent truncation. The original fifteen logits are retained during expansion so that a different batching context cannot perturb the baseline. Hashes of the corpus, queries, and judgments link the saved inputs to the recorded runs.

\subsection{Development decisions and follow-up status}
For SciFact, the fixed development sample comprises one fifth of the official training queries; NFCorpus and FiQA use their supplied development splits. The base search considers
\begin{equation}
\begin{aligned}
\beta&\in\{0,0.25,0.5\},\\
\lambda&\in\{0,0.01,0.025,0.05\}.
\end{aligned}
\label{eq:development_grids}
\end{equation}
All stochastic settings use seeds zero through four. Eligibility requires mean nDCG@5 and Recall@8 to match or exceed L6-only on the same shortlist within $10^{-12}$. Among eligible settings, the lexicographic order first maximizes nDCG@5, then Recall@8; it then minimizes redundancy, followed by $\lambda$ and finally $\beta$.

SciFact selects a zero diversity weight, which the rule permits. At that setting the relevance-ranked base incumbent already optimizes the unquantized relevance objective. Calling the search GAS does not by itself imply a nonzero diversity effect; the objective, frozen weights, and search procedure have to be interpreted separately.

The L6 MMR grid is 0, 0.1, 0.25, 0.5, 0.75, and 1: six values. The recorded relevance-first preference selects zero diversity weight on these collections. Its measured row represents one operating point, not the complete MMR tradeoff frontier. Cosine-guided controls retain their own frozen settings. Legacy threshold selection considers $\{0.3,0.5,0.7\}$ using development nDCG@5 and the declared tie preferences.

The deterministic 100-query confirmation samples exclude historical development, pilot, and test IDs. The descriptive gate compares refinement with its matched base: mean nDCG@5 must be equal, the mean-recall ratio at least 0.98, and mean redundancy strictly lower. NFCorpus lost 2.1561\% relative recall in confirmation and failed the gate. A recorded decision continued to the full test without changing the method or global gate; the later outcome does not turn the confirmation into a pass.
% TODO(author): Explain how the recall-retention threshold was chosen, if a rationale is documented.

Earlier test outcomes had been observed during development. Although the reported parameters were not retuned against the newer follow-up judgments, those official tests are not an untouched holdout for the redesign. The journal diagnostics retrospectively describe the same saved runs.

\subsection{Controls and comparison fairness}
The original evaluation has twenty methods, with seven expanded-pool policies added for protected-refinement runs. The controls include dense cosine, dense/BM25 RRF with rank constant 60, native cosine and L6 MMR, cosine-guided GAS, L6-only ranking, base L6-guided GAS, greedy and exact optimizers, random search, legacy threshold Grover, and L2/L6 neural models. Cross-encoders on dense fifteen and dense sixty-four use different pools from L6-only on hybrid fifteen.

The seven expanded-pool policies comprise constrained GAS, constrained local-swap greedy, exact float, exact quantized, matched random search, global union-L6 ranking, and union-L6 with the protected prefix. Each constrained policy uses the actual baseline prefix for the seed and the same tail feasibility rules. The random refiner draws $\max(1,2Q)$ uniform subsets in addition to its initial incumbent, with $Q$ denoting the seed's recorded refinement phase count. That matches a declared objective-access allowance, not every operation of the compact simulator.

The exact controls independently enumerate their float or quantized objectives. The GAS controller receives neither their optimum nor the marked-state count. A comparison across pools describes complete policies; one among constrained refiners with identical features isolates their search outcomes for the same objective and admissible domain. These comparisons answer different questions.

\subsection{Metrics, inference, and hardware}
We use nDCG@5, Recall@8, and the redundancy measure in~\eqref{eq:redundancy} as principal metrics; nDCG@8 shows effects at ranks six through eight. Discounted gains use linear judgment values and logarithmic rank discount~\cite{jarvelin2002gain}, and saved rankings are checked with NIST's trec\_eval~\cite{trec_eval}. For a five-seed policy, we first average within each query and then across queries. Seeds do not create independent relevance judgments for a query.

For protected-run inference, the recorded procedure uses 10,000 paired bootstrap resamples with seed zero and 10,000 paired sign-randomization draws with seed one. Differences no larger than $10^{-12}$ are ties. Holm correction covers nine comparators by four metrics per protected dataset. The earlier FiQA nDCG audit covers six declared comparators and removes numerical averaging artifacts before testing. We retain those families and their pointwise intervals rather than choosing a new family to change apparent significance.

The recorded host runs Linux and Python 3.12.3 on an Intel Core i7-1355U with 12 logical CPUs and approximately 15.3\,GiB of RAM. Compute libraries use two threads. Recorded software versions include NumPy 1.26.4, SciPy 1.13.1, CPU PyTorch 2.12.0, Transformers 4.57.6, Sentence Transformers 3.4.1, Qiskit 0.46.0, and Aer 0.13.3~\cite{javadi2024qiskit}. Online latency excludes model loading, corpus/cache restoration, and artificial-query warm-up. Session and host-load records are available, but the experiments were not run in a dedicated latency laboratory.

Online time includes the applicable retrieval, lexical, neural, and selection stages; reused upstream work is charged to each policy that requires it. Base and refinement times are summed for each seeded request before combined selector percentiles are calculated. Request P95 is not divided by the number of seeds. Section~\ref{sec:diagnostics} decomposes mean cost additively, without treating sums of component percentiles as pipeline percentiles.
